\section{Declaración de uso de IA generativa}

\paragraph{Declaración.} Se usó la herramienta Claude (Anthropic) como apoyo en tres tareas: explorar con código la hoja principal del archivo de la ANT (conteos, consistencia y faltantes), comparar alternativas de definición y apoyar la redacción del informe. También se usó para generar los modelos UML y verificar las cifras con el procedimiento del Anexo~B. El grupo revisó cada afirmación y es responsable de la exactitud del contenido; las cifras de la sección~6 pueden reproducirse con ese procedimiento.

\begin{table}[H]
\centering
\small
\caption{Recomendaciones de IA verificadas, corregidas o descartadas}
\label{tab:bitacora}
\begin{tabularx}{\textwidth}{p{0.7cm}X X p{3.0cm}}
\toprule
\# & Recomendación de la IA & Verificación & Resultado \\
\midrule
1 & Excluir las columnas de condición de los participantes porque revelan el desenlace & El conteo de condición \texttt{FALLECIDO} coincide con \texttt{FALLECIDOS} en los 17.236 registros (Tabla~\ref{tab:a4}) & Verificada y adoptada \\
2 & Tratar la causa probable como un conjunto de información separado (B) y no como una variable más & Se contrastó con la retroalimentación recibida y con las proporciones por causa, que van del 1,9\,\% al 23,1\,\% (Tabla~\ref{tab:a5}) & Corregida respecto a la propuesta inicial \\
3 & Describir el conjunto B y las ventanas móviles en términos de entrenar y comparar modelos & En el análisis del problema no se entrena nada; esos elementos son condiciones de evaluación & Corregida; se describen en la sección~9 y en las Tablas~\ref{tab:a6} y~\ref{tab:a7} \\
\bottomrule
\end{tabularx}
\end{table}
